% Plantilla «presentación» — creada por el wizard de LaTeX Studio (M1-F4).
\documentclass[aspectratio=169]{beamer}

\usetheme{Madrid}
\usecolortheme{default}

\usepackage[T1]{fontenc}
\usepackage[utf8]{inputenc}
\usepackage[spanish]{babel}

\title{Presentación}
\subtitle{Subtítulo}
\author{Nombre Apellido}
\institute{Institución}
\date{\today}

\begin{document}

\begin{frame}
  \titlepage
\end{frame}

\begin{frame}{Agenda}
  \tableofcontents
\end{frame}

\section{Introducción}

\begin{frame}{Introducción}
  \begin{itemize}
    \item Primer punto: contexto del tema.
    \item Segundo punto: objetivo de la presentación.
  \end{itemize}
\end{frame}

\section{Desarrollo}

\begin{frame}{Desarrollo}
  Contenido principal de la presentación.
\end{frame}

\begin{frame}{Conclusiones}
  \begin{itemize}
    \item Resumen de ideas clave.
    \item Próximos pasos.
  \end{itemize}
\end{frame}

\end{document}
